\documentclass{article}

% NeurIPS 2024 style
\usepackage[preprint]{neurips_2024}

\usepackage[utf8]{inputenc}
\usepackage[T1]{fontenc}
\usepackage{hyperref}
\usepackage{url}
\usepackage{booktabs}
\usepackage{amsfonts}
\usepackage{amsmath}
\usepackage{amssymb}
\usepackage{nicefrac}
\usepackage{microtype}
\usepackage{xcolor}
\usepackage{graphicx}
\usepackage{subcaption}
\usepackage{multirow}
\usepackage{array}
\usepackage{enumitem}
\usepackage{algorithm}
\usepackage{algorithmicx}
\usepackage{algpseudocode}

% Convenience
\newcommand{\eg}{\textit{e.g.}}
\newcommand{\ie}{\textit{i.e.}}
\newcommand{\etal}{\textit{et al.}}
\newcommand{\relppl}{\text{rel-PPL}}
\newcommand{\kv}{\mathrm{KV}}
\newcommand{\kvpatch}{\texttt{kvpatch}}
\definecolor{darkgreen}{rgb}{0.0,0.5,0.0}
\definecolor{darkred}{rgb}{0.7,0.0,0.0}

\title{Subspace PolarQuant: KV Cache Compression\\via PCA Projection and Rotation Quantization}

\author{%
  Corey Petty \\
  Institute of Free Technology \\
  \texttt{corbin@corpetty.com}
}

\begin{document}

\maketitle

%------------------------------------------------------------------------------
\begin{abstract}
We study compression of the key-value (KV) cache in transformer language models
via projection into a learned PCA subspace followed by PolarQuant
rotation-quantization. Through 21 experiments across five architectures
(Qwen3-1.7B/14B/32B, Mistral-7B, Phi-4, Llama-3.1-8B), we establish that
\textbf{truncation error from dimension reduction categorically dominates
quantization noise}: $k{=}64$, 16-bit (pure truncation) produces $2.48{\times}$
perplexity degradation while $k{=}128$, 4-bit (zero truncation, aggressive
quantization) produces only $1.05{\times}$. This motivates a compression target of
$k/d_\text{head} \geq 0.875$, yielding $4.27{\times}$ compression at $14\%$ perplexity
cost on Qwen3-14B. We further find that \textbf{V compression is not viable at
any $k < d_\text{head}$ across all tested architectures}---including Llama-3.1
which lacks QK-norm---indicating that V vectors are intrinsically
high-dimensional, not an artifact of a specific normalization scheme. K
compression generalizes across architectures ($4{\times}$ compression, $<5\%$ PPL on
Llama-3.1; $<20\%$ PPL on all tested models at $k{=}112$), context lengths (stable
to 40K tokens), and calibration domains. We release \kvpatch{}, a library
enabling drop-in KV compression via \texttt{patch(model,~tokenizer,~k=112)}.
\end{abstract}

%------------------------------------------------------------------------------
\section{Introduction}
\label{sec:intro}

Inference with large language models is increasingly memory-bound at long context
lengths. For a 14-billion-parameter model like Qwen3-14B at 32K tokens, the KV
cache alone requires approximately 10\,GB of GPU memory---comparable to the model
weights themselves. This limits the practical context window on consumer hardware
and increases serving costs proportionally with sequence length.

KV cache compression aims to reduce this memory footprint at inference time
without modifying model weights. The central challenge is that KV vectors must be
reconstructed with sufficient fidelity to preserve attention score distributions,
which determines the model's output quality.

We investigate a method combining two techniques: \textbf{PCA subspace projection},
which projects each $d_\text{head}$-dimensional KV vector into a $k$-dimensional
principal subspace fitted on calibration data, and \textbf{PolarQuant}
\cite{han2025polarquant}, a learned-rotation quantization scheme that applies an
orthogonal rotation before scalar quantization to decorrelate dimensions.

Our primary contribution is an empirical characterization of this method across
scales and architectures, producing three actionable findings:

\begin{enumerate}[leftmargin=1.5em, itemsep=2pt]
  \item \textbf{Truncation error is the binding constraint.} The critical design
    variable is $k$ (subspace dimension), not bit depth. Practitioners should
    maximize $k$ before minimizing bits.
  \item \textbf{V compression is universally hard.} Key vectors compress well
    across all tested architectures. Value vectors resist subspace compression at
    any $k < d_\text{head}$, regardless of whether the architecture uses QK-norm.
    This is an intrinsic property of V variance structure.
  \item \textbf{K compression is architecture-agnostic and long-context-stable.}
    K subspace compression generalizes across Qwen3, Mistral, Phi-3, and
    Llama-3.1, and maintains stable quality through 40K-token contexts with
    offline calibration.
\end{enumerate}

%------------------------------------------------------------------------------
\section{Background and Related Work}
\label{sec:background}

\paragraph{KV cache compression.}
Prior work falls broadly into three categories: \emph{token eviction}
\cite{zhang2023h2o, ge2023model}, which discards KV entries for less important
tokens; \emph{quantization} \cite{hooper2024kvquant, liu2024kvsharer}, which
reduces bit precision of stored vectors; and \emph{dimensionality reduction},
which projects KV vectors to lower-rank representations. GQA
\cite{ainslie2023gqa} reduces KV head count at training time; our work instead
reduces the per-head vector dimensionality at inference time without retraining.

\paragraph{PolarQuant.}
PolarQuant \cite{han2025polarquant} addresses the non-uniform distribution of
transformer activations by applying a learned orthogonal rotation $R$ before
quantization, spreading variance more uniformly across dimensions. We use
PolarQuant as the quantizer within our projected subspace.

\paragraph{PCA subspace structure.}
The KV vectors produced by transformer attention heads exhibit low effective rank.
In our analysis of Qwen3-14B, K vectors have a mean effective rank (90\% variance
threshold) of 29.6 out of $d_\text{head}{=}128$, while V vectors have mean
effective rank 54.2. This asymmetry foreshadows the finding that V compression is
harder.

%------------------------------------------------------------------------------
\section{Method}
\label{sec:method}

\subsection{Notation}
\label{sec:notation}

Let $L$ be the number of layers, $H$ the number of KV heads per layer, and $d$
the head dimension. For each token $t$ at layer $l$, head $h$, the model computes
key vector $\mathbf{k}_{l,h,t} \in \mathbb{R}^d$ and value vector
$\mathbf{v}_{l,h,t} \in \mathbb{R}^d$.

\subsection{Calibration}
\label{sec:calibration}

Given a calibration corpus $\mathcal{C}$ ($\sim$2K tokens of generic text), we
collect all key vectors via forward hooks. For each $(l, h)$ pair we compute:
\begin{equation}
  \boldsymbol{\mu}_{l,h} = \frac{1}{|\mathcal{C}|}\sum_t \mathbf{k}_{l,h,t},
  \qquad
  \mathbf{U}_{l,h} \in \mathbb{R}^{d\times d}
  \text{ from SVD of centered key matrix.}
\end{equation}

Calibration cost: one forward pass on $\sim$2K tokens. Basis storage:
$L \times H \times d^2 \times 4\,\text{bytes} \approx 45\,\text{MB}$ for
Qwen3-14B ($d=128$).

\subsection{Compression and Reconstruction}
\label{sec:compression}

\begin{algorithm}[t]
  \caption{Subspace PolarQuant --- Compress / Reconstruct}
  \label{alg:compress}
  \begin{algorithmic}[1]
    \Require key $\mathbf{k} \in \mathbb{R}^d$, basis
      $\mathbf{U}_{l,h}$, mean $\boldsymbol{\mu}_{l,h}$,
      rotation $R\in\mathbb{R}^{k\times k}$, integers $k$, $n_\text{bits}$
    \State $\tilde{\mathbf{k}} \gets \mathbf{k} - \boldsymbol{\mu}_{l,h}$
      \Comment{Center}
    \State $\mathbf{z} \gets \mathbf{U}_{l,h}[:, {:}k]^\top \tilde{\mathbf{k}}$
      \Comment{Project to top-$k$ subspace}
    \State $\mathbf{z}' \gets R\,\mathbf{z}$
      \Comment{Rotate (PolarQuant)}
    \State $q \gets \mathrm{Quantize}(\mathbf{z}', n_\text{bits})$
      \Comment{Uniform scalar quantization}
    \State \textbf{store} $(q, \text{scale}, \text{offset})$
      \Comment{$k \cdot n_\text{bits}$ bits vs $d \cdot 16$ bits}
    \Statex
    \State $\mathbf{z}' \gets \mathrm{Dequantize}(q)$
    \State $\mathbf{z} \gets R^\top \mathbf{z}'$
    \State $\hat{\mathbf{k}} \gets \mathbf{U}_{l,h}[:, {:}k]\,\mathbf{z}
           + \boldsymbol{\mu}_{l,h}$
      \Comment{Reconstruct in $\mathbb{R}^d$}
  \end{algorithmic}
\end{algorithm}

The compression ratio is:
\begin{equation}
  \text{CR} = \frac{d \cdot 16}{k \cdot n_\text{bits}}.
  \label{eq:cr}
\end{equation}
For $k{=}112$, $n{=}4$, $d{=}128$: $\text{CR} = 2048/448 = 4.57{\times}$.

\subsection{V Treatment}
\label{sec:v_treatment}

Our experiments establish that V compression fails at $k < d$ for all tested
architectures (Section~\ref{sec:v_failure}). The recommended deployment therefore
uses:
\begin{itemize}[itemsep=2pt]
  \item \textbf{K}: Subspace PolarQuant ($k{=}112$, 4-bit) --- $4.57{\times}$ compression
  \item \textbf{V}: Full-dimensional PolarQuant ($k{=}d{=}128$, 4-bit) --- $4.00{\times}$
  \item \textbf{Combined KV}: $\approx 4.27{\times}$
\end{itemize}

%------------------------------------------------------------------------------
\section{Experiments}
\label{sec:experiments}

\paragraph{Protocol.} All experiments use three evaluation passages (biology,
computer science, history) with perplexity as the primary metric. Relative PPL
$= \text{PPL}_\text{compressed} / \text{PPL}_\text{baseline}$. Viability
threshold: $\relppl < 1.20$ (20\% degradation budget). All experiments run on a
single NVIDIA RTX 3090 (24\,GB) via SLURM.

%--------------------------------------
\subsection{Truncation vs.\ Quantization Error}
\label{sec:truncation_vs_quant}

\paragraph{Setup.} We sweep $k \in \{64, 96, 112, 128\}$ and bits $\in \{4, 6,
8, 16\}$ on Qwen3-14B-AWQ (40 layers, 8 KV heads, $d_\text{head}{=}128$).

\begin{table}[h]
  \centering
  \caption{Relative PPL for $k \times$ bits sweep on Qwen3-14B-AWQ. Values
  below $1.20$ are highlighted (\textcolor{darkgreen}{green}).}
  \label{tab:bitrate_sweep}
  \begin{tabular}{@{}lcccc@{}}
    \toprule
    $k$ & 4-bit & 6-bit & 8-bit & 16-bit \\
    \midrule
    64  & 3.11 & 2.46 & 2.43 & 2.43 \\
    96  & 1.25 & 1.17 & 1.17 & 1.17 \\
    112 & \textcolor{darkgreen}{1.14} & \textcolor{darkgreen}{1.09} & \textcolor{darkgreen}{1.08} & \textcolor{darkgreen}{1.08} \\
    128 & \textcolor{darkgreen}{1.05} & \textcolor{darkgreen}{1.00} & \textcolor{darkgreen}{1.00} & --- \\
    \bottomrule
  \end{tabular}
\end{table}

Moving from $k{=}64$ to $k{=}128$ at fixed 4-bit reduces $\relppl$ from
$3.11{\times}$ to $1.05{\times}$ ($\Delta{=}2.06{\times}$). Moving from 4-bit to 16-bit at
fixed $k{=}64$ reduces $\relppl$ from $3.11{\times}$ to $2.43{\times}$
($\Delta{=}0.68{\times}$). \textbf{Truncation error is $\sim$3$\times$ more
impactful than quantization error} on the scale of this sweep, and the gap
widens at larger $k$ where quantization noise is already small.

%--------------------------------------
\subsection{End-to-End Perplexity at Scale}
\label{sec:scale}

\paragraph{Qwen3 family.}
\begin{table}[h]
  \centering
  \caption{K compression ($k{=}112$/4-bit) across Qwen3 model sizes. CR =
  $4.27{\times}$ for all.}
  \label{tab:qwen_scale}
  \begin{tabular}{@{}llcc@{}}
    \toprule
    Model & Params & Rel PPL & Min viable $k$ \\
    \midrule
    Qwen3-1.7B    & 1.7B  & 1.32$\times$ & 128 \\
    Qwen3-14B-AWQ & 14B   & \textcolor{darkgreen}{1.14$\times$} & 112 \\
    Qwen3-32B-AWQ & 32B   & \textcolor{darkgreen}{1.06$\times$} & 96  \\
    \bottomrule
  \end{tabular}
\end{table}

Compression tolerance scales with model size. Qwen3-32B tolerates $k{=}96$ at
$1.09{\times}$ PPL ($4.57{\times}$ CR), consistent with the over-parameterization
hypothesis: larger models develop lower-rank functional KV subspaces.

\paragraph{Cross-architecture comparison.}
\begin{table}[h]
  \centering
  \caption{K compression ($k{=}112$/4-bit) across architectures.
  $^\dagger$K-only (V excluded; see Section~\ref{sec:v_failure}).}
  \label{tab:cross_arch}
  \begin{tabular}{@{}llccc@{}}
    \toprule
    Model & Architecture & Rel PPL & Min viable $k$ & CR at min-$k$ \\
    \midrule
    Mistral-7B-v0.3  & Mistral  & \textcolor{darkgreen}{1.07$\times$} & 64  & $5.33{\times}$ \\
    Phi-4-AWQ        & Phi-3    & \textcolor{darkgreen}{1.10$\times$} & 64  & $5.33{\times}$ \\
    Qwen3-14B-AWQ    & Qwen3    & \textcolor{darkgreen}{1.14$\times$} & 112 & $4.27{\times}$ \\
    Llama-3.1-8B$^\dagger$ & LLaMA & \textcolor{darkgreen}{1.04$\times$} & 112 & $4.57{\times}$ \\
    \bottomrule
  \end{tabular}
\end{table}

Mistral and Phi-3 tolerate $k{=}64$ ($50\%$ truncation) within the budget.
Llama-3.1 achieves the best K-only result ($1.04{\times}$) of any tested
architecture.

%--------------------------------------
\subsection{Long-Context Stability}
\label{sec:long_context}

\begin{table}[h]
  \centering
  \caption{Relative PPL vs.\ context length ($k{=}128$/4-bit and $k{=}96$/4-bit
  on War and Peace).}
  \label{tab:long_ctx}
  \begin{tabular}{@{}lcccccccc@{}}
    \toprule
    Config & 512 & 1K & 2K & 4K & 8K & 16K & 32K & 40K \\
    \midrule
    $k{=}128$/4-bit
      & \textcolor{darkgreen}{1.11} & \textcolor{darkgreen}{1.11} & \textcolor{darkgreen}{1.09} & \textcolor{darkgreen}{1.06} & \textcolor{darkgreen}{1.05} & \textcolor{darkgreen}{1.06} & \textcolor{darkgreen}{1.06} & \textcolor{darkgreen}{1.09} \\
    $k{=}112$/4-bit
      & 1.68 & 1.47 & 1.54 & 1.39 & 1.34 & 1.42 & 1.65 & 1.67 \\
    $k{=}96$/4-bit
      & 2.30 & 2.10 & 1.91 & 1.62 & 1.66 & 1.71 & 1.83 & 1.85 \\
    $k{=}64$/4-bit
      & 14.99 & 10.78 & 7.18 & 6.00 & 5.19 & 4.59 & 4.37 & 4.26 \\
    \bottomrule
  \end{tabular}
\end{table}

$k{=}128$/4-bit is stable within $1.05$--$1.11{\times}$ across all context lengths
with no accumulation. Per-position analysis confirms compression errors are
uniformly distributed across token positions.

\textbf{Basis drift.} PCA basis overlap between early (tokens 0--2K) and late
(tokens 35K--40K) document positions: K overlap $= 0.825$, V overlap $= 0.702$.
V drifts more than K but stabilizes. The offline basis remains representative at
40K tokens.

%--------------------------------------
\subsection{Task Performance: Needle-in-Haystack Retrieval}
\label{sec:needle}

\begin{table}[h]
  \centering
  \caption{Retrieval accuracy for facts inserted at varied depths and context
  lengths. Config: $k{=}128$/4-bit vs.\ baseline.}
  \label{tab:needle}
  \begin{tabular}{@{}lcccccc@{}}
    \toprule
    Config & 4K & 8K & 16K & 32K & Overall \\
    \midrule
    Baseline           & 93\% & 93\% & 87\%  & 100\% & 93\% \\
    $k{=}128$/4-bit    & 93\% & 93\% & 100\% & 100\% & \textbf{97\%} \\
    $k{=}96$/4-bit     & 100\% & 93\% & 100\% & 27\%  & 80\% \\
    \bottomrule
  \end{tabular}
\end{table}

$k{=}128$/4-bit matches or exceeds baseline accuracy at all context lengths,
achieving $97\%$ overall vs.\ $93\%$ baseline. The marginal improvement at 16K is
consistent with a mild regularization effect from compression noise.
$k{=}96$/4-bit collapses at 32K (27\%), consistent with its PPL instability above
16K context.

%--------------------------------------
\subsection{V Compression: Architecture-Independent Failure}
\label{sec:v_failure}

\paragraph{Hypothesis.}
Initial experiments on Qwen3-14B showed V compression failing at all $k < 128$.
We hypothesized this might be a Qwen3-specific artifact of QK-norm (RMSNorm
applied to $\mathbf{q}$/$\mathbf{k}$ outputs), which may force K into a
lower-dimensional manifold while leaving V structure undistorted.

\paragraph{Experiment: Llama-3.1-8B-Instruct-AWQ.}
Llama-3.1 uses standard GQA without QK-norm, enabling a direct test.

\begin{table}[h]
  \centering
  \caption{V compression results on Llama-3.1-8B (no QK-norm). QK-norm
  hypothesis rejected.}
  \label{tab:v_failure}
  \begin{tabular}{@{}llcc@{}}
    \toprule
    Config & Description & Rel PPL & Viable? \\
    \midrule
    K-only $k{=}112$/4-bit  & K compressed, V full rank        & \textcolor{darkgreen}{1.042$\times$} & \checkmark \\
    V-only $k{=}112$/4-bit  & V compressed, K full rank        & 12.14$\times$ & $\times$ \\
    K+V $k{=}128$/4-bit     & Both full-rank 4-bit             & \textcolor{darkgreen}{1.085$\times$} & \checkmark \\
    K+V $k{=}124$/4-bit     & 3.1\% truncation of V            & 3.45$\times$ & $\times$ \\
    \bottomrule
  \end{tabular}
\end{table}

The QK-norm hypothesis is \textbf{rejected}. Llama-3.1 without QK-norm shows
identical V compression failure---$k_V{=}112$ produces $12.14{\times}$ PPL,
\emph{worse} than Qwen3 at the same $k$. The hard cliff at $k_V{=}124$
($3.45{\times}$ PPL) is architecture-independent: the $\sim$30\% of V variance
residing in tail principal components (dimensions 113--128 for $d_\text{head}{=}128$)
is load-bearing signal.

%--------------------------------------
\subsection{Adaptive K-Scheduling}
\label{sec:adaptive}

\paragraph{Layer sensitivity (Exp.~16).}
Ablating each of 40 layers at $k{=}64$/4-bit reveals strong heterogeneity. Layer
37 is most sensitive ($+1.93$ PPL delta). Layers 2, 5, 20, 25, 27, 29, 38 show
\emph{negative} sensitivity (slight improvement under compression, consistent
with mild regularization). Late layers near the LM head are most sensitive.

\paragraph{Rank-proportional scheduling (Exp.~18).}
Assigning $k_l$ proportional to each layer's effective rank, subject to a
mean-$k$ budget constraint, outperforms uniform $k$:

\begin{table}[h]
  \centering
  \caption{Adaptive K-scheduling vs.\ uniform at matched mean-$k$ budgets.}
  \label{tab:adaptive}
  \begin{tabular}{@{}lccc@{}}
    \toprule
    Budget (mean $k$) & Uniform rel PPL & Rank-prop.\ rel PPL & $\Delta$ \\
    \midrule
    110 & 1.153$\times$ & \textcolor{darkgreen}{1.132$\times$} & $-1.8\%$ \\
    \bottomrule
  \end{tabular}
\end{table}

The rank-proportional policy achieves mean $k{=}110$ (vs.\ uniform $k{=}112$) with
\emph{lower} PPL---$1.8\%$ better quality at $1.8\%$ less memory. This policy is
derived directly from the calibration forward pass at zero additional cost.

%--------------------------------------
\subsection{Cross-Domain Calibration}
\label{sec:cross_domain}

Calibrating on one domain (fiction, code, news, dialogue) and evaluating on
another incurs a modest penalty at $k{=}128$/4-bit: worst cross-domain pair is
within $0.05{\times}$ $\relppl$ of same-domain calibration. At $k{=}96$/4-bit,
sensitivity increases: code$\to$news gives $\relppl{=}2.70{\times}$ vs.\ universal
calibration at $1.63{\times}$.

\textbf{Practical recommendation:} A single calibration pass on a diverse 2K-token
corpus generalizes safely at $k{=}128$. For aggressive $k{=}96$, universal
calibration (mixed-domain) is required.

%------------------------------------------------------------------------------
\section{Analysis: Why Does Truncation Dominate?}
\label{sec:analysis}

The PCA decomposition partitions a KV vector's variance into orthogonal components
ordered by magnitude. The tail components---while small in $\ell_2$
norm---represent systematic directions that are always absent under truncation.

\paragraph{Attention score sensitivity.} The attention score for query $\mathbf{q}$
and key $\mathbf{k}$ is $\text{sim}(\mathbf{q}, \mathbf{k}) = \mathbf{q}^\top
\mathbf{k}$. Compression error in $\mathbf{k}$ propagates directly to attention
scores. For rare but important tokens (names, numbers, technical terms), attention
must be precise; the tail components of $\mathbf{k}$ may be systematically
aligned with the query directions that select such tokens.

\paragraph{Layer-wise error accumulation.} In a 40-layer transformer, each
layer's compressed KV output becomes input to the next layer's attention.
Quantization noise averages out across layers (random, uncorrelated). Truncation
error accumulates: the same dimensions are always absent. This explains the
disproportionate performance drop at $k{=}64$ in end-to-end evaluation versus
single-layer distortion metrics.

\paragraph{V intrinsic dimensionality.} The $\sim$30\% of V variance in tail
dimensions (113--128) is qualitatively different from K tail variance. Online
basis updating (Exp.~19) and larger subspaces (Exp.~20) both fail to recover V
quality, suggesting the signal is structurally spread across all dimensions rather
than localized in a low-dimensional manifold. This is consistent with V vectors
representing token content (arbitrary), while K vectors represent positional and
relational structure (more constrained).

%------------------------------------------------------------------------------
\section{The \texttt{kvpatch} Library}
\label{sec:kvpatch}

\subsection{API}
\label{sec:api}

\begin{verbatim}
from kvpatch import patch, calibrate, unpatch

# One-call: auto-calibrate and patch
patch(model, tokenizer, k=112, bits=4)

# Explicit calibration for basis reuse
basis = calibrate(model, tokenizer, k=112, bits=4,
                  save_path="basis.pkl")
patch(model, basis=basis)

# Remove compression
unpatch(model)
\end{verbatim}

\subsection{Architecture Support}
\label{sec:arch_support}

\kvpatch{} auto-detects architectures by inspecting the model config and module
structure: Qwen2/Qwen3 (AWQ and standard), LLaMA/Mistral, Phi-3, Falcon, and a
generic fallback scanning \texttt{named\_modules} for \texttt{k\_proj} and
\texttt{v\_proj}. The current hook-based implementation adds $\sim$13$\times$
decode overhead due to CPU roundtrip; a CUDA-fused backend would reduce this to
$\sim$1.2$\times$ estimated (see Section~\ref{sec:limitations}).

\subsection{Memory Impact}
\label{sec:memory}

At $k{=}112$, 4-bit, for Qwen3-14B at 32K context:
\begin{itemize}[itemsep=0pt]
  \item Uncompressed K cache: 5.24\,GB
  \item Compressed K cache: 1.15\,GB (\textbf{4.09\,GB saved})
\end{itemize}
This enables running 32K context on a single 24\,GB GPU that would otherwise
require two GPUs or reduced context length.

%------------------------------------------------------------------------------
\section{Discussion}
\label{sec:discussion}

\paragraph{What this means for practitioners.}
When tuning compression, increase $k$ first. Going from $k{=}112$ to $k{=}128$
improves PPL more than doubling bit depth from 4 to 8 at the same $k$. K and V
are different problems: K compresses to $4.57{\times}$ with modest quality loss; V
does not benefit from subspace compression and should use full-dimensional
quantization at $4.00{\times}$. Offline calibration on 2K tokens of generic text
is sufficient for cross-domain, cross-context generalization.

\paragraph{Large models compress better.}
Within the Qwen3 family, compression tolerance increases monotonically with scale.
This is consistent with the over-parameterization hypothesis: larger models
develop lower-rank functional KV subspaces because they have more redundant
capacity.

\subsection{Limitations}
\label{sec:limitations}

\paragraph{Hook overhead.} The current implementation uses PyTorch forward hooks
with CPU-side operations, incurring $\sim$13$\times$ decode slowdown. A fused
CUDA kernel for project-rotate-quantize is required for production use.

\paragraph{Small-model limitation.} Models below $\sim$5B parameters (Qwen3-1.7B
in our tests) require $k{=}d$ to stay within the 20\% PPL budget, capturing no
benefit from subspace projection. The method is most impactful for $\geq$7B
models.

\paragraph{Basis storage.} Per-(layer, head) PCA bases require $\sim$45\,MB for
Qwen3-14B. Modest but non-zero; scales with model depth and KV head count.

\subsection{Future Work}
\label{sec:future}

\begin{itemize}[itemsep=2pt]
  \item \textbf{Fused CUDA kernel.} Would reduce hook overhead to $\sim$1.2$\times$
    estimated, making the method throughput-positive at batch size $\geq$4.
  \item \textbf{100K+ context scaling.} V basis drift (overlap 0.702 at 40K)
    likely worsens at 100K+. Periodic basis refresh without the quality-neutral
    overhead of online updates (Exp.~19) may be necessary.
  \item \textbf{V-norm at training time.} If QK-norm reduces K intrinsic
    dimensionality, an analogous V-norm might make V compression viable. This
    hypothesis is testable via fine-tuning on a V-normalized model variant.
  \item \textbf{Non-standard KV formats.} MLA \cite{deepseek2024} and other
    latent KV compression schemes may interact with subspace quantization in
    unexplored ways.
\end{itemize}

%------------------------------------------------------------------------------
\section{Conclusion}
\label{sec:conclusion}

We have systematically characterized subspace PolarQuant KV compression across 21
experiments on six model variants. The central finding---that truncation error
categorically dominates quantization noise---resolves a longstanding ambiguity in
KV cache compression design: maximize subspace dimension $k$ before minimizing
bit depth. At $k{=}112$/4-bit, Qwen3-14B achieves $4.27{\times}$ KV cache compression
with $14\%$ perplexity cost, stable across 40K context lengths, calibrated on two
thousand tokens of generic text, and applicable via three lines of Python code.

The secondary finding---that V compression is architecture-independently
intractable below full rank---closes the most natural extension of this method and
redirects future work toward improving V vector compressibility, possibly at
training time via V-norm analogues to QK-norm.

%------------------------------------------------------------------------------
\bibliography{refs}
\bibliographystyle{plainnat}

\end{document}
